\section{Software: Arcane Framework}
\label{sec:WP7:Arcane Framework:software}

\begin{table}[h!]
    \centering
    { \setlength{\parindent}{0pt}
    \def\arraystretch{1.25}
    \arrayrulecolor{numpexgray}
    {\fontsize{9}{11}\selectfont
    \begin{tabular}{!{\color{numpexgray}\vrule}p{.4\textwidth}!{\color{numpexgray}\vrule}p{.6\textwidth}!{\color{numpexgray}\vrule}}
        \rowcolor{numpexgray}{\rule{0pt}{2.5ex}\color{white}\bf Field} & {\rule{0pt}{2.5ex}\color{white}\bf Details} \\
        \rowcolor{white}\textbf{Consortium} & \begin{tabular}{l}
CEA\\
IFPEN\\
\end{tabular} \\
        \rowcolor{numpexlightergray}\textbf{Exa-MA Partners} & \begin{tabular}{l}
CEA\\
\end{tabular} \\
        \rowcolor{white}\textbf{Contact Emails} & \begin{tabular}{l}
lydie.grospellier@cea.fr\\
\end{tabular} \\
        \rowcolor{numpexlightergray}\textbf{Supported Architectures} & \begin{tabular}{l}
CPU or GPU\\
\end{tabular} \\
        \rowcolor{white}\textbf{Repository} & \href{https://github.com/arcaneframework/framework}{https://github.com/arcaneframework/framework} \\
        \rowcolor{numpexlightergray}\textbf{License} & \begin{tabular}{l}
OSS: apache-2\\
\end{tabular} \\
        \bottomrule
    \end{tabular}
    }}
    \caption{WP7: Arcane Framework Information}
\end{table}

\subsection{Software Overview}
\label{sec:WP7:Arcane Framework:summary}

Arcane is a production-grade multi-physics simulation framework developed by CEA and IFPEN. Within WP7 it acts as a reference
codebase to validate the shared build and benchmarking methodology: Arcane exercises the complete pipeline from reproducible
container builds (Spack- and Docker-based) to CI-driven performance tests and energy measurements. The framework is also used
to stress WP7's data-management guidelines thanks to its rich mesh and field I/O requirements.

\begin{table}[h!]
    \centering
    { 
        \setlength{\parindent}{0pt}
        \def\arraystretch{1.25}
        \arrayrulecolor{numpexgray}
        {
            \fontsize{9}{11}\selectfont
            \begin{tabular}{!{\color{numpexgray}\vrule}p{.25\linewidth}!{\color{numpexgray}\vrule}p{.6885\linewidth}!{\color{numpexgray}\vrule}}
    
    \rowcolor{numpexgray}{\rule{0pt}{2.5ex}\color{white}\bf Features} &  {\rule{0pt}{2.5ex}\color{white}\bf Short Description }\\ 
    
\rowcolor{white}    Environment & C++17 framework with WP7-maintained Spack recipes and container images validated on Jean Zay and IFPEN clusters. \\
\rowcolor{numpexlightergray}    Energy & Integration of the WP7 power-profiling plug-in (RAPL/NVML) within Arcane's accelerator abstraction layer. \\
\rowcolor{white}    Physics & Multi-physics finite-volume modules covering CFD, combustion, and neutronics demonstrators. \\
\end{tabular}
        }
    }
    \caption{WP7: Arcane Framework Features}
\end{table}


\subsection{Parallel Capabilities}
\label{sec:WP7:Arcane Framework:performances}


Arcane combines distributed-memory parallelism (MPI) with on-node acceleration through OpenMP and CUDA/HIP back-ends exposed by
its \texttt{Arcane::Accelerator} layer. Production runs typically target tens of thousands of CPU cores or multi-GPU partitions
on the CEA TGCC and Jean Zay systems. Hybrid assemblies---for instance MPI+OpenMP on CPU partitions and MPI+CUDA on GPU nodes---
are configured via a unified runtime that also manages mesh repartitioning and halo exchanges. The framework interfaces with
third-party solvers such as PETSc and HYPRE, which are packaged through WP7's shared Spack environment to guarantee binary
compatibility across applications.


\subsection{Initial Performance Metrics}
\label{sec:WP7:Arcane Framework:metrics}

The first WP7 benchmarking campaign focuses on a weak-scaling CFD case (mixing tank) instrumented with Extrae and the WP7 energy
plug-in. Inputs and reference outputs are archived on the WP7 Zenodo community (deposition pending publication) and mirrored in the Exa-MA object store to
comply with FAIR principles. On 256 GPU-enabled ranks the application sustains 78\% parallel efficiency while keeping
energy-to-solution within 5\% of the CPU-only baseline thanks to asynchronous I/O pipelines. The primary bottleneck identified so
far is mesh redistribution during adaptive refinement; mitigation work explores lightweight metadata exchanges and tighter
coupling with the ParMmg workflow maintained under WP1.

\subsubsection{Benchmark \#1}
\begin{itemize}
    \item \textbf{Description:} Mixing-tank CFD case with 120M cells, targeting the Jean Zay GPU partition (4\,×\,A100 per node) to qualify weak scaling and energy efficiency.
    \item \textbf{Benchmarking Tools Used:} Extrae and the WP7 energy plug-in record MPI/accelerator traces, while ReFrame orchestrates runs and collects wall-clock, memory, and joule metrics.
    \item \textbf{Input/Output Dataset Description:}
        \begin{itemize}
            \item \textbf{Input Data:} Initial mesh and boundary condition set (8~GB HDF5 archive) available through the WP7 Zenodo community archive.
            \item \textbf{Output Data:} Time-series of velocity and temperature fields (compressed HDF5) plus profiling logs stored alongside the DOI artefact.
            \item \textbf{Data Repository:} Zenodo collection ``Exa-MA WP7 -- Arcane benchmarks''.
        \end{itemize}
    \item \textbf{Results Summary:} 0.92\,s/step at 256 ranks (78\% parallel efficiency versus 32 ranks), energy-to-solution reduced by 18\% compared with CPU-only run.
    \item \textbf{Challenges Identified:} Dynamic load-balancing triggers additional mesh exchanges; optimisation work is tracking reductions in metadata traffic and improved GPU staging.
\end{itemize}

\subsection{12-Month Roadmap}
\label{sec:WP7:Arcane Framework:roadmap}

\begin{itemize}
    \item \textbf{Data improvements:} Release a curated Zenodo collection with automated metadata extraction from ReFrame runs and long-term retention on the TGCC object store (Q1~2026).
    \item \textbf{Methodology application:} Extend the WP7 CI pipeline to cover mixed-precision solver variants and nightly energy regressions on Jean Zay pre-production nodes (Q2~2026).
    \item \textbf{Results retention:} Publish a living benchmark dashboard (Grafana) fed by WP7 metrics exports and cross-referenced from the Exa-MA knowledge base (Q3~2026).
\end{itemize}